% Bagian: sets-functions-relations
% Bab: arithmetization
% Subbagian: refleksi

\documentclass[../../../include/open-logic-section]{subfiles}

\begin{document}

\olfileid[id]{sfr}{arith}{ref}
\olsection{Beberapa Refleksi Filosofis}

Demikianlah rincian teknisnya. Namun, apa yang telah dicapai oleh semua itu?

Hampir tidak terbantahkan bahwa semua itu memberi kita matematika murni yang
{indah}. Selain itu, terdapat beberapa pencapaian konseptual yang
mendalam. Menyadari bahwa Sifat Kelengkapan mengungkapkan perbedaan penting
antara bilangan real dan bilangan rasional merupakan wawasan yang mendalam.
Selain itu, konstruksi eksplisit bilangan real sebagai potongan Dedekind
memberi landasan yang kokoh bagi pokok bahasan analisis. Kita mengetahui bahwa
gagasan \emph{lapangan terurut lengkap} bersifat koheren, karena
potongan-potongan tersebut membentuk tepat lapangan semacam itu.

Meskipun demikian, kita perlu mengemukakan beberapa keberatan terhadap
pencapaian-pencapaian tersebut.

Pertama, tidak jelas bahwa memikirkan bilangan real dalam kerangka potongan
\emph{lebih} ketat secara matematis daripada memikirkan bilangan real dalam kerangka ekspansi
desimalnya yang sudah dikenal (yang mungkin tak berhingga). ``Konstruksi''
bilangan real yang disebut terakhir ini agak menyerupai konstruksi bilangan
real melalui barisan Cauchy; tetapi sebenarnya konstruksi itu pada intinya
telah dikenal para matematikawan sejak awal abad ke-17 (lihat
\olref[cauchy]{sec}). Peningkatan ketelitian matematis yang sebenarnya berasal dari
kesadaran bahwa bilangan real memiliki Sifat Kelengkapan; kemampuan membangun
bilangan real sebagai himpunan tertentu mungkin, dengan sendirinya, tidak
begitu menarik.

Lebih tidak jelas lagi bahwa aritmetisasi bilangan bulat, atau bilangan
rasional, yang jauh lebih mudah, meningkatkan ketelitian matematis dalam ranah-ranah
tersebut. Di sini ada baiknya kita membuat suatu pengamatan sederhana. Setelah
\emph{membangun} bilangan bulat sebagai kelas-kelas ekuivalensi dari pasangan
terurut bilangan asli, lalu membangun bilangan rasional sebagai kelas-kelas
ekuivalensi dari pasangan terurut bilangan bulat, dan kemudian membangun
bilangan real sebagai himpunan bilangan rasional, kita segera
\emph{melupakan} konstruksi-konstruksi tersebut. Secara khusus: tidak seorang
pun akan pernah ingin \emph{menggunakan} konstruksi-konstruksi tersebut dalam
suatu pembuktian matematis (kecuali, tentu saja, pembuktian bahwa
konstruksi-konstruksi itu berperilaku sebagaimana yang dimaksudkan). Jauh lebih
mudah berbicara langsung tentang suatu bilangan real daripada berbicara
tentang suatu himpunan dari himpunan dari himpunan dari himpunan dari himpunan
dari himpunan dari himpunan bilangan asli.
%(Sebab bilangan real dibangun sebagai himpunan bilangan rasional; sedangkan bilangan rasional dibangun sebagai kelas-kelas ekuivalensi dari pasangan terurut bilangan bulat, yakni himpunan dari himpunan dari himpunan bilangan bulat; dan seterusnya.)
%2/3 = [2,3]
%	yakni {<2,3>, ... }
% 	{ {{2},{2,3}}, ... }
%
%bilangan real adalah
%	himpunan bilangan rasional
%
%bilangan rasional adalah
%	himpunan dari himpunan dari himpunan bilangan bulat
%
%bilangan bulat adalah
%	himpunan dari himpunan dari himpunan bilangan asli
	
Yang paling meragukan dari semuanya ialah bahwa definisi-definisi ini memberi
tahu kita apa \emph{hakikat} bilangan bulat, bilangan rasional, atau bilangan
real \emph{secara metafisis}. Artinya, patut diragukan bahwa bilangan real
(misalnya) \emph{adalah} himpunan tertentu (dari himpunan dari
himpunan\ldots). Penghalang utama bagi pandangan semacam itu ialah bahwa
konstruksinya dapat dilakukan dengan banyak cara berbeda. Dalam kasus
bilangan real, ada beberapa konstruksi berbeda yang benar-benar menarik (lihat
\olref[cauchy]{sec}). Namun, berikut cara yang sungguh sepele untuk memperoleh
beberapa konstruksi berbeda: seperti dalam \olref[sfr][rel][ref]{sec}, kita
dapat mendefinisikan pasangan terurut secara agak berbeda; seandainya kita
menggunakan gagasan alternatif tentang pasangan terurut ini, konstruksi kita
akan berhasil sama baiknya dengan sebelumnya, tetapi kita akan berakhir
dengan objek-objek berbeda. Dengan demikian, terdapat banyak konstruksi
berbasis teori himpunan yang saling bersaing bagi bilangan bulat, bilangan rasional,
dan bilangan real. Maka, akan semena-mena (dan memalukan) jika kita mengklaim
bahwa bilangan bulat (misalnya) adalah \emph{himpunan-himpunan ini}, dan bukan
\emph{himpunan-himpunan itu}. (Seperti dalam \olref[sfr][rel][ref]{sec}, ini
merupakan suatu contoh argumen yang dimasyhurkan oleh
\citealt{Benacerraf1965}.)

Ada satu hal lebih lanjut yang patut dikemukakan: ada sesuatu yang cukup
\emph{ganjil} pada konstruksi kita. Kita mulai dengan bilangan asli. Kemudian
kita membangun bilangan bulat, dan membangun ``$0$ dari bilangan bulat'',
yakni $ \equivrep{0,0}{\Intequiv}$. Akan tetapi, $0 \neq
\equivrep{0,0}{\Intequiv}$. Bahkan, berdasarkan konstruksi kita,
\emph{tidak ada} bilangan asli yang merupakan bilangan bulat. Namun, hal ini
tampak sangat berlawanan dengan intuisi. Bahkan, dalam
\olref[sfr][set][imp]{sec}, tanpa banyak argumentasi kita mengklaim bahwa
$\Nat \subseteq \Rat$. Jika konstruksi tersebut memberi tahu kita secara
persis \emph{hakikat} bilangan, klaim ini jelas keliru.

Jika kita melihatnya secara menyeluruh, apa kesimpulan kita? Dengan bekerja dalam
teori himpunan naif, dan memanfaatkan bilangan asli, kita dapat
\emph{memperlakukan} bilangan bulat, bilangan rasional, dan bilangan real
sebagai himpunan-himpunan tertentu. Dalam arti itu, kita dapat
\emph{menanamkan} teori-teori tentang entitas tersebut di dalam suatu teori
himpunan. Namun, makna filosofis penanaman ini sama sekali tidak sederhana.

Tentu saja, semua ini bukan kata akhir!{} Intinya hanya ini. Untuk menunjukkan
bahwa aritmetisasi bilangan real \emph{memang} memiliki arti filosofis yang
mendalam, diperlukan suatu argumen \emph{filosofis} tambahan.

%Selain itu: sepanjang bab ini, kita telah memperlakukan
%bilangan asli seolah-olah sekadar \emph{diberikan} kepada kita. Namun, apa
%yang dapat kita lakukan dengannya? Itulah pokok bahasan bab berikutnya.

\end{document}
